\section{Related Work}
\label{sec:related}

\subsection{Collaborative Filtering}

Classical collaborative filtering (CF) methods recommend items based on the overlap of user interaction histories. User--user and item--item approaches compute cosine or Pearson similarity on the interaction matrix; matrix factorisation decomposes it into low-rank user and item latent factors, often optimised with Bayesian Personalised Ranking~\cite{rendle2009bpr} to handle implicit feedback. Neural Collaborative Filtering (NCF)~\cite{he2017ncf} replaces the inner product with an MLP, learning non-linear user--item interaction patterns. While these methods achieve competitive ranking quality, they are agnostic to content: two users who happen to share popular tracks appear similar regardless of the musical reasons behind those choices, amplifying the popularity bias inherent in implicit feedback data. Our work uses KNN-CF and a popularity baseline to contextualise graph-based results against these established approaches.

\subsection{Content-Based and Hybrid Music Recommendation}

Content-based approaches characterise tracks via audio or symbolic features and recommend items similar to a user's listened catalogue. In the symbolic domain, jSymbolic~\cite{mckay2018jsymbolic} and music21~\cite{music21docs} extract hundreds of descriptors from MIDI files spanning pitch, rhythm, harmony, and instrumentation. The challenge is high dimensionality and redundancy among correlated features; autoencoders have been used to compress such representations before recommendation. Hybrid models fuse content representations with collaborative signals: our XGBoost baseline retrieves candidates via autoencoder embeddings and re-ranks with a gradient-boosted LTR model, isolating the contribution of graph structure by providing a strong content-plus-behaviour baseline without relational context.

\subsection{Knowledge Graph-Enhanced Recommendation}

Knowledge graphs encode structured domain knowledge that standard CF cannot exploit. KGCN~\cite{wang2019kgcn} performs graph convolution on KG neighbours to propagate entity semantics into item representations. KGAT~\cite{wang2019kgat} extends this with attention over the KG, weighting relations by their relevance to user preferences end-to-end. LightGCN~\cite{he2020lightgcn} strips graph convolution to its essential propagation step, showing that learnable weights and non-linearities are unnecessary for homogeneous user--item graphs. For music specifically, Liu et al.~\cite{liu2024kgmusic} combine KG structural features with multi-task feature learning, achieving gains over CF-only baselines by incorporating genre and artist entity representations.

Our approach differs from these in three ways: (i)~we construct the KG from scratch using established music vocabularies and Wikidata alignment rather than mapping to a pre-built KG; (ii)~we use KG embeddings (RotatE, ComplEx) as node feature initialisers for a \textit{heterogeneous} GNN rather than aggregating KG neighbourhoods in a separate module; and (iii)~we incorporate 1,525-dimensional symbolic audio features via an autoencoder, providing musical content that KG structure alone cannot represent.

\subsection{Heterogeneous Graph Neural Networks}

Standard GNNs aggregate over all neighbours identically. In a music KG, however, a user's connection to a genre node carries fundamentally different semantics from a track's connection to a key node, and conflating these through shared weights loses this type information. The Heterogeneous Graph Transformer (HGT)~\cite{hu2020hgt} addresses this by learning separate transformation and attention parameters for each (source-type, edge-type, target-type) triple, preserving the graph's typed structure throughout message passing. This makes HGT particularly well suited to our multi-relational music KG, where node types span users, tracks, artists, genres, instruments, keys, modes, tempo classes, and decades. An additional advantage over opaque embedding models is that HGT's attention weights can be traced back to specific edge types and neighbouring entities, providing recommendation rationales that are directly grounded in KG structure — an important property for a domain where users benefit from understanding why a track was suggested.

\subsection{Knowledge Graph Embeddings for Structural Initialisation}

Transductive KG embedding models (KGE) learn dense vectors for all entities from relational triples alone, without requiring any external features. RotatE~\cite{sun2019rotate} and ComplEx~\cite{trouillon2016complex} are complex-valued models well suited to asymmetric and compositional relation patterns. In many KG-enhanced recommenders, KGE vectors are used as pre-computed item side features fed into a separate recommendation model. Our use differs: we train KGE on the non-interaction triples of our KG and use the resulting vectors to initialise \emph{all} entity node types in the heterogeneous GNN, including structural nodes (genres, keys, instruments) that have no other numeric representation. This separates structural pre-training from interaction modelling and prevents target-signal leakage into the node features.
